\chapter{Ejemplos Prácticos de Control y Telemetría}
\label{ch:practical-examples}

\section{Introducción}
\label{sec:practical-intro}

Este capítulo proporciona ejemplos prácticos que demuestran cómo interactúan la aplicación Android
y el firmware \textbf{MICSBOL/ESP32\_BT\_Controller} en un robot real o un sistema embebido.

Del lado Android se transmite continuamente un \textbf{estado de control} (sticks, perillas e interruptores),
y también se transmiten \textbf{eventos de una sola vez} cuando el usuario presiona botones.

Del lado del ESP32:
\begin{itemize}
    \item decodifica los paquetes entrantes por Bluetooth Serial (SPP),
    \item aplica los valores de control recibidos a salidas de hardware reales (PWM, pines de dirección, interruptores),
    \item y opcionalmente envía telemetría de vuelta a la app Android (paneles, indicadores, gráficas).
\end{itemize}

\section{Datos de Control desde la Aplicación Android}
\label{sec:control-data}

\subsection{Paquete de estado RC (\texttt{rcStatePacket})}
El firmware almacena el estado de control más reciente en la estructura definida en \texttt{packets.h}:

\begin{verbatim}
    rcStatePacket
\end{verbatim}

Contiene los siguientes campos:

\begin{itemize}
    \item \texttt{leftStickX}, \texttt{leftStickY}
    \item \texttt{rightStickX}, \texttt{rightStickY}
    \item \texttt{leftKnob}, \texttt{rightKnob}
    \item \texttt{switches} (máscara de bits)
\end{itemize}

En el código del repositorio, los canales de sticks y perillas se tratan como valores de 12 bits y se mapean en el
rango \textbf{0..4095} (ver \texttt{control.cpp}, \texttt{hal\_output.cpp}).

\subsection{Paquete de evento (\texttt{rcEventPacket})}
Los botones se transmiten como pequeños paquetes de evento, y el ID del evento se pasa a:

\begin{verbatim}
    controlHandleEvent(byte eventId)
\end{verbatim}

Los eventos son útiles para acciones que deben ser \emph{momentáneas} (p.\ ej., disparar, reset, bocina) en lugar de un
valor analógico continuo.

\section{Cómo el ESP32 Recibe Paquetes}
\label{sec:receiver}

Los bytes Bluetooth entrantes se procesan en:

\begin{verbatim}
    receiver.cpp  ->  handleBluetooth()
\end{verbatim}

El receptor mantiene un buffer circular y busca encabezados de paquetes.
En la implementación actual del firmware:

\begin{itemize}
    \item Encabezado RC State: \texttt{0xAA 0x55}
    \item Encabezado Event: \texttt{0xBB 0x66}
\end{itemize}

Una vez decodificado un evento, \texttt{controlHandleEvent()} se llama inmediatamente.

\section{Dónde se Aplican los Valores de Control}
\label{sec:control-layer}

El bucle principal de control se implementa en:

\begin{verbatim}
    control.cpp  ->  controlUpdate()
\end{verbatim}

Esta función lee valores de \texttt{rcStatePacket} y los envía a la capa de abstracción de hardware
(\texttt{hal\_output.cpp}):

\begin{itemize}
    \item PWM del servo de dirección
    \item comandos de motor izquierdo/derecho
    \item PWM de paneo de cámara
    \item brillo de LED e intensidad de zumbador (PWM)
    \item salidas de interruptores (máscara de bits decodificada en salidas por canal)
\end{itemize}

\section{Ejemplo 1: Control de Dirección Usando el Joystick Izquierdo}
\label{sec:example-steering}

El firmware mapea el eje horizontal del joystick izquierdo a un valor PWM tipo servo:

\begin{verbatim}
    rcStatePacket.data.leftStickX
\end{verbatim}

En \texttt{control.cpp}, la dirección se genera con:

\begin{verbatim}
    static uint32_t mapServo(uint16_t v) {
        return map(v, 0, 4095, 205, 410);
    }

    uint32_t steerPWM = mapServo(rcStatePacket.data.leftStickX);
    halSetSteering(steerPWM);
\end{verbatim}

\subsection*{Interpretación práctica}
\begin{itemize}
    \item El joystick Android produce un valor de 0 a 4095.
    \item El ESP32 lo mapea a un rango PWM más pequeño (205..410) adecuado para su configuración PWM.
\end{itemize}

\noindent Esto se usa comúnmente para:
\begin{itemize}
    \item dirección de auto RC
    \item dirección de ruedas delanteras de robot
    \item un timón accionado por servo
\end{itemize}

\section{Ejemplo 2: Control de Velocidad de Motores en Tracción Diferencial}
\label{sec:example-motors}

El firmware usa los ejes verticales de los joysticks como canales de motor izquierdo/derecho:

\begin{verbatim}
    rcStatePacket.data.leftStickY
    rcStatePacket.data.rightStickY
\end{verbatim}

En \texttt{control.cpp}:

\begin{verbatim}
    uint32_t motorL = rcStatePacket.data.leftStickY;
    uint32_t motorR = rcStatePacket.data.rightStickY;

    halSetMotorLeft(motorL);
    halSetMotorRight(motorR);
\end{verbatim}

\subsection{Mapeo de dirección + velocidad (comportamiento implementado)}
En \texttt{hal\_output.cpp}, cada valor de motor se interpreta con una banda muerta (dead-band):

\begin{itemize}
    \item \textbf{0..2000:} pin de dirección LOW (reversa) y PWM mapeado a 0..4095
    \item \textbf{2001..2094:} banda muerta, objetivo del motor se establece en 0
    \item \textbf{2095..4095:} pin de dirección HIGH (adelante) y PWM mapeado a 0..4095
\end{itemize}

Este diseño reduce el jitter cuando el joystick está cerca del centro.

\subsection{Rampa suave (comportamiento avanzado)}
En lugar de escribir PWM instantáneamente, el PWM del motor hace una rampa hacia un valor objetivo. Una interrupción por timer actualiza
las salidas de motor cada 20 ms, moviendo la velocidad por un paso fijo (ver \texttt{updateMotors()} en
\texttt{hal\_output.cpp}). Esto mejora la manejabilidad y reduce picos de corriente.

\section{Ejemplo 3: Control de Paneo de Cámara}
\label{sec:example-camera}

El eje horizontal del joystick derecho se usa para paneo de cámara:

\begin{verbatim}
    uint32_t panPWM = mapServo(rcStatePacket.data.rightStickX);
    halSetCameraPan(panPWM);
\end{verbatim}

\noindent Usos típicos en la vida real:
\begin{itemize}
    \item cámara de rover de inspección
    \item posicionamiento de cámara FPV
    \item montaje de sensor tipo torreta
\end{itemize}

\section{Ejemplo 4: Usar Perillas para Controlar PWM de LED y Zumbador}
\label{sec:example-knobs}

La app Android proporciona dos perillas, transmitidas como valores 0..4095:

\begin{verbatim}
    rcStatePacket.data.leftKnob
    rcStatePacket.data.rightKnob
\end{verbatim}

En \texttt{control.cpp}, estas se mapean directamente a salidas PWM mediante el HAL:

\begin{verbatim}
    halSetLed(rcStatePacket.data.leftKnob);
    halSetBuzzer(rcStatePacket.data.rightKnob);
\end{verbatim}

\subsection*{Usos prácticos}
\begin{itemize}
    \item brillo de LED (faros, underglow, luz de trabajo)
    \item volumen / intensidad del zumbador (alertas, pitidos de confirmación)
    \item reutilizar perillas para escalado de límite de motor o trim de servo (modificación personalizada del firmware)
\end{itemize}

\section{Ejemplo 5: Controlar Salidas Digitales con Interruptores}
\label{sec:example-switches}

Los interruptores se transmiten como una máscara de bits en un solo byte:

\begin{verbatim}
    byte sw = rcStatePacket.data.switches;
\end{verbatim}

En \texttt{control.cpp}, el firmware itera a través de los índices de switch:

\begin{verbatim}
    for (int i = 0; i < 6; i++) {
        halSetSwitch(i, (sw >> i) & 1);
    }
\end{verbatim}

\subsection*{Diferencia importante por modo}
La cantidad de switches soportados depende del modo de proyecto (ver \texttt{hal\_output.cpp}):

\begin{itemize}
    \item \textbf{FULL\_RC\_MODE:} solo los índices 0..3 se escriben a pines GPIO (4 salidas de switch).
    \item \textbf{FULL\_RC\_MODE\_MCP:} los índices 0..5 se escriben usando salidas del MCP23017 (6 salidas de switch).
\end{itemize}

\noindent Usos en el mundo real:
\begin{itemize}
    \item switch de habilitar/armar (seguridad)
    \item faros / baliza / underglow
    \item salidas de relé (bomba, solenoide, alimentación de accesorio)
    \item selección de modo (manual vs asistido)
\end{itemize}

\section{Ejemplo 6: Manejo de Eventos de Botón (Salidas de Pulso)}
\label{sec:example-events}

Los eventos se manejan mediante:

\begin{verbatim}
    controlHandleEvent(byte eventId)
\end{verbatim}

En la implementación del repositorio, los eventos llaman a \texttt{halPulseSwitch(index)}. El mapeo de índices difiere
según el modo de proyecto. Por ejemplo, en \texttt{FULL\_RC\_MODE}:

\begin{verbatim}
    case 0x01: halPulseSwitch(0); break;
    case 0x02: halPulseSwitch(1); break;
    case 0x03: halPulseSwitch(0); break;
    case 0x04: halPulseSwitch(1); break;
\end{verbatim}

\subsection{Pulsos no bloqueantes}
\texttt{halPulseSwitch()} usa \texttt{pulseStart()} que genera un \textbf{pulso no bloqueante}
(por defecto 50 ms) y requiere que \texttt{pulseUpdate()} se llame en el lazo principal.

\noindent Acciones típicas de evento:
\begin{itemize}
    \item bocina / efecto de sonido corto
    \item disparador ``tomar foto'' para una cámara
    \item resetear un subsistema (re-zero de IMU, línea de reset de controlador)
    \item relé momentáneo (pestillo de puerta, golpe de actuador)
\end{itemize}

\section{Datos de Telemetría Enviados a la Aplicación Android}
\label{sec:telemetry-data}

La telemetría es generada por:

\begin{verbatim}
    telemetry.cpp
    telemetry_source.cpp
\end{verbatim}

El diseño está intencionalmente dividido:
\begin{itemize}
    \item \textbf{\texttt{telemetry.cpp}} construye y envía paquetes (solo escritura).
    \item \textbf{\texttt{telemetry\_source.cpp}} provee valores (lecturas de hardware o valores inyectados por debug).
\end{itemize}

\subsection{Tipos de telemetría y tasas (implementado)}
\begin{itemize}
    \item \textbf{Indicadores:} cada 500 ms (\(\approx 2\,Hz\)).
    \item \textbf{Gráficas (plots):} cada 50 ms (\(\approx 20\,Hz\)).
    \item \textbf{Paneles:} retransmitidos brevemente cuando los valores cambian (ventana de reenvío para robustez).
\end{itemize}

\subsection{Config-on-connect}
En una nueva conexión, el ESP32 envía un \textbf{paquete de configuración} que incluye nombres de gráficas/paneles/indicadores
(p.\ ej., \texttt{Volts}, \texttt{Amps}, \texttt{RPMs}, \texttt{Left}, \texttt{Right}, \texttt{Throttle}, \texttt{Battery}).
Esto permite que la app Android etiquete elementos de UI automáticamente.

\section{Ejemplo 7: Telemetría de Batería como Indicador}
\label{sec:example-battery}

El paquete de indicadores del firmware incluye:
\begin{itemize}
    \item \texttt{analogValue} (0--100)
    \item \texttt{batteryLevel} (0--100)
    \item \texttt{digitalMask} (bitfield)
\end{itemize}

En modo normal, la batería se lee desde un pin ADC y se mapea a un porcentaje (ver \texttt{hwReadIndicatorBattery()}
en \texttt{input\_hw.cpp}):

\begin{verbatim}
    return map(analogRead(PIN_ANALOG_IND2), 0, 4095, 0, 100);
\end{verbatim}

\noindent Nota práctica de cableado: el voltaje de batería debe escalarse con un divisor de voltaje para que el ADC del ESP32
nunca supere 3.3V.

\section{Ejemplo 8: Paneles de Telemetría desde Entradas Analógicas}
\label{sec:example-panels}

La telemetría de panel usa dos valores de 16 bits (\texttt{leftPanelValue}, \texttt{rightPanelValue}). La implementación de hardware
por defecto lee valores analógicos y los escala de 0..4095 a 0..65535:

\begin{verbatim}
    uint16_t hwReadNumeric1() {
        uint32_t v = analogRead(PIN_NUMERIC1);
        return (uint16_t)((v * 65535UL) / 4095);
    }
\end{verbatim}

Esto es adecuado para:
\begin{itemize}
    \item sensores de distancia con salida analógica
    \item realimentación con potenciómetro
    \item señales de sensor escaladas (después de acondicionamiento)
\end{itemize}

\section{Ejemplo 9: Señales de Gráfica (0..255 Gráficas Rápidas)}
\label{sec:example-plots}

Las gráficas transmiten tres valores de 8 bits (0..255) a \(\approx 20\,Hz\). La implementación de hardware por defecto lee
valores ADC y los mapea a 0..255:

\begin{verbatim}
    uint8_t hwReadPlot1() {
         return map(analogRead(PIN_NUMERIC1), 0, 4095, 0, 255);
         }
    uint8_t hwReadPlot2() { \\
        return map(analogRead(PIN_NUMERIC2), 0, 4095, 0, 255);
        }
    uint8_t hwReadPlot3() {
        return map(analogRead(PIN_ANALOG_IND2), 0, 4095, 0, 255);
        }
\end{verbatim}

\noindent Ideas prácticas para gráficas:
\begin{itemize}
    \item tendencia de voltaje de batería (caída bajo aceleración)
    \item tendencia de sensor de corriente (detectar atascos)
    \item tendencia de velocidad derivada de RPM/encoder
\end{itemize}

\section{Ejemplo 10: Indicadores de Máscara Digital}
\label{sec:example-digital-mask}

La telemetría de indicadores incluye un byte de máscara digital que puede representar hasta 8 señales on/off.
En \texttt{input\_hw.cpp}, se construye una máscara a partir de lecturas GPIO de indicadores (o desde entradas del MCP23017 en modo MCP).

Esto es útil para mostrar de forma compacta:
\begin{itemize}
    \item finales de carrera
    \item banderas de falla
    \item sensores de parachoques
    \item señales de ``armado'' / ``listo''
\end{itemize}

\section{Resumen}
\label{sec:practical-summary}

El sistema soporta comunicación bidireccional:

\begin{itemize}
    \item \textbf{Android \(\rightarrow\) ESP32:} estado RC continuo + eventos discretos.
    \item \textbf{ESP32 \(\rightarrow\) Android:} paneles, indicadores y gráficas con configuración automática de etiquetas.
\end{itemize}

Al combinar entradas de control (sticks, perillas, switches, eventos) con retroalimentación de telemetría en vivo, puedes construir
robots y sistemas embebidos que sean \textbf{fáciles de operar} y \textbf{fáciles de depurar en tiempo real}.